% ECONOMICS_PROBLEM: {"id": "C72-488", "title": "Critical Bernoulli densities for expected-polynomial Nash computation", "jel_code": "C72", "source_id": "OP-0138"}
% !TeX program = xelatex
\documentclass[11pt,a4paper]{article}
\usepackage[margin=23mm]{geometry}
\usepackage{fontspec}
\setmainfont{Latin Modern Roman}
\usepackage{amsmath,amssymb,mathtools}
\usepackage{xcolor,enumitem,booktabs,longtable,array,xurl,hyperref}
\definecolor{muted}{HTML}{53636C}
\hypersetup{colorlinks=true,urlcolor=blue,linkcolor=blue}
\setlength{\parindent}{0pt}
\setlength{\parskip}{5pt}
\newcommand{\R}{\mathbb R}
\newcommand{\E}{\mathbb E}
\newcommand{\N}{\mathbb N}
\newcommand{\ind}{\mathbf 1}
\newcommand{\Var}{\operatorname{Var}}
\newcommand{\Cov}{\operatorname{Cov}}
\newcommand{\supp}{\operatorname{supp}}
\DeclareMathOperator*{\argmin}{arg\,min}
\DeclareMathOperator*{\argmax}{arg\,max}
\newcommand{\entrymeta}[3]{{\sffamily\small\color{muted}\textbf{\texttt{#1}}\quad|\quad #2\par #3\par}}
\newcommand{\Problemref}[1]{\texttt{#1}}
\newcommand{\JELref}[2]{\texttt{#2} (JEL \texttt{#1})}
\begin{document}
\section*{Critical Bernoulli densities for expected-polynomial Nash computation}
\textbf{Problem ID:} \texttt{C72-488}\quad \textbf{JEL:} \texttt{C72}\par
% BEGIN ECONOMICS STATEMENT
\label{problem:OP-0138}
% GAME_THEORY_PROBLEM: {"local_id": "GT-008", "original_number": 8, "kind": "P", "entry_kind": "statement_only", "published": false, "review_required": true, "category": "game_theory"}
\label{game:GT-008}
{\sffamily\small\color{muted}
\textbf{GT-008} | Average-case equilibrium complexity\\
\textbf{P} | Explicit remaining parameter regimes in the cited source.
\par}
\subsection*{Definitions and mathematical statement}
For $n\geq1$ and $p_n\in[0,1]$, sample $A,B\in\{0,1\}^{n\times n}$ so that their $2n^2$ entries are mutually independent Bernoulli variables with parameter $p_n$. Exact Nash equilibrium means the two best-response inequalities in GT-005. For each fixed $c>0$ and each of the exponents $a\in\{1/2,1\}$, put $p_n=\min\{1,cn^{-a}\}$. Determine whether an always-correct exact-equilibrium algorithm exists with
\[
 \mathbb E_{A,B,\xi}[T_{\mathcal A}(A,B;\xi)]\leq C_{a,c}n^{d_{a,c}}
 \qquad\text{for every }n\geq1.
\]
The regimes not covered by the source's displayed sufficient bounds include
\[
 a=\tfrac12,\quad e^{-52}2^{-8}<c<0.977;
 \qquad a=1,\quad c>0.3849\text{ fixed}.
\]
The exact asymptotic version may replace $p_n=cn^{-a}$ by a specified sequence with $p_n/(cn^{-a})\to1$; endpoint claims must then follow the source's precise inequalities.
\subsection*{Short English statement}
Close the remaining density gaps in average-case exact equilibrium algorithms for random win--lose games.
\subsection*{Variants and scope boundaries}
The source also covers $p_n\geq(\log n)^9/n$ in a sparse-density regime. That bound is not $c/n$ for a fixed constant $c$, and must not be used to claim that the fixed-$c$ critical line has been solved. High-probability existence of a small-support equilibrium alone does not control expected fallback time.
\subsection*{Source relationship and status}
This is source-date residual scope, not an exhaustive audit of every subsequent random-game paper. The threshold constants are sufficient conditions established in [S1], not conjectured sharp transitions.
\subsection*{References}
\begingroup\footnotesize\raggedright
{\footnotesize\raggedright
\textbf{[S1]} Andrea Collevecchio, G\'abor Lugosi, Adrian Vetta, and Rui-Ray Zhang (2025). \emph{Finding a Nash equilibrium of a random win-lose game in expected polynomial time}. arXiv:2510.12846. Locator: Abstract, Table 1, Theorem 1, and Sections 5--6. \url{https://arxiv.org/pdf/2510.12846}\par}
\par\endgroup

\subsection*{Common conventions: Source mathematical conventions}

\textbf{Finite games.} For a finite set $A$, $\Delta(A)$ is its probability
simplex. Players choose independent mixed strategies in Nash equilibrium;
correlation is allowed only when explicitly part of the solution concept.
In a game with payoff functions $u_i$ and strategy profile $x$, an additive
$\varepsilon$ Nash equilibrium satisfies
\[
 u_i(x)\geq u_i(x_i',x_{-i})-\varepsilon
 \quad\text{for every player $i$ and every admissible deviation $x_i'$}.
\]
For numerical approximation, payoffs are normalized as locally specified.
An $\varepsilon$ payoff residual is distinct from distance at most
$\varepsilon$ to the set of exact equilibria. Point convergence, convergence
of empirical play, and convergence in distance to an equilibrium set are
also distinct.

\textbf{Computational representations.} Unless stated otherwise, a complexity
question takes rational data with binary encoding; polynomial time is measured
in total bit length. A fixed-accuracy polynomial algorithm is not necessarily
polynomial in $1/\varepsilon$, and neither is necessarily polynomial in
$\log(1/\varepsilon)$. Succinct games and value, demand, or sampling oracles
must declare their interfaces. Exact real solutions require an explicit
representation; an unspecified list of arbitrary real numbers is not a
finite computational output. A claimed algorithmic barrier is meaningful
only for its specified model and reductions.

\textbf{Dynamic games.} Histories, information, observation, and allowable
randomization are part of the model. A stationary strategy depends only on
the current state; a positional strategy in a deterministic graph game
chooses one action at each controlled vertex. Nash equilibrium at an initial
state and equilibrium after every history are different requirements.
An expectation of a pathwise $\liminf$ payoff, a $\liminf$ of expected averages,
a limiting expected average, and a uniform finite-horizon equilibrium are
different criteria. None is silently substituted for another.

\textbf{Allocation and mechanisms.} A complete allocation assigns every item
exactly once unless otherwise stated. Zero-valued goods, negative values,
capacity constraints, feasibility, lotteries, and copies of goods affect
fairness definitions and are specified locally. Dominant-strategy incentive
compatibility, Bayesian incentive compatibility, universal truthfulness,
and truthfulness in expectation impose different inequalities. Whenever
an objective ratio has zero denominator, its convention is stated or those
instances are excluded.

\textbf{Infinite-dimensional models.} Expectations, integrals, stochastic
equations, and optimization objectives require their local regularity,
integrability, filtrations, and admissible domains. A backward--forward
mean-field system is a boundary-value problem; it is not an initial-value
semiflow without an additional construction. Long-time, large-population,
and vanishing-noise limits require an order of limits and a convergence
topology. If a minimum need not be attained, the objective is its infimum
or supremum and the approximation requirement is stated separately.
% END ECONOMICS STATEMENT
\end{document}
